\section{Log canonical pairs and continuation of the program}\label{sec:lc}

We continue to work over $\mathbb C$.  We first deduce termination for
elementary $\mathbb Q$-factorial dlt programs from
Theorem~\ref{thm:klt-termination}.  We then lift a prescribed small tail of
an arbitrary lc program to such a program.  The lift is an auxiliary
construction in the proof: it does not replace the given initial model
in Definition~\ref{def:program}.
For the standard definition of dlt pairs, see \cite{KM98}.

\subsection{Elementary dlt steps and special termination}

An elementary step in this subsection means the usual divisorial
contraction or flip of a negative extremal ray on a
$\mathbb Q$-factorial dlt pair.  Their existence follows from the klt
contraction and flip theorems and relative base point freeness.

\begin{lemma}\label{lem:dlt-elementary}
Let $(U,\Delta)$ be a projective $\mathbb Q$-factorial dlt pair with rational boundary,
and let $R$ be a $(K_U+\Delta)$-negative extremal ray.  Its elementary
contraction exists.  If it is birational and small, its flip exists,
is projective, and is a $(K_U+\Delta)$-flip.  Divisorial contractions and
flips preserve $\mathbb Q$-factoriality and the dlt property.
\end{lemma}

\begin{proof}
Put $D=K_U+\Delta$ and $F=\lfloor\Delta\rfloor$.  For a sufficiently
small positive rational number $\varepsilon$, the pair
\[
  (U,\Delta_\varepsilon),\qquad
  \Delta_\varepsilon=\Delta-\varepsilon F,
\]
is klt, and $D_\varepsilon=K_U+\Delta_\varepsilon$ is negative on $R$.
Here $F$ is $\mathbb Q$-Cartier by $\mathbb Q$-factoriality.  Decreasing
an effective $\mathbb Q$-Cartier part of a dlt boundary preserves dlt,
and a dlt boundary with no coefficient one is klt; see
\cite[Corollary~2.39 and Proposition~2.41]{KM98}.  Negativity on $R$
is an open condition on $\varepsilon$, since it can be checked on one
nonzero numerical class in this ray.

The klt cone and contraction theorem gives the contraction
$f:U\to V$ of $R$, with normal projective target and connected fibres.
It has relative Picard number one.  If $f$ is small, klt flip existence
gives a projective small map $f^+:U^+\to V$ on which
$D_\varepsilon^+=K_{U^+}+\Delta_\varepsilon^+$ is relatively ample;
see \cite[Corollary~1.4.1]{BCHM}.  The output is
$\mathbb Q$-factorial; the elementary statements are also recorded in
\cite[Lemma~3.10.2]{BCHM}.

We check the positive sign for the full boundary.  In $N^1(U/V)$ there
is a positive rational number $t$ such that
\[
 D\equiv_V tD_\varepsilon.
\]
Indeed both classes have negative degree on $R$, this relative space
has dimension one, and their degrees on an integral curve are
rational.  The $\mathbb Q$-Cartier divisor
$L=D-tD_\varepsilon$ is numerically trivial over $V$.  A Cartier
multiple of $L$ is relatively nef, and its difference from
$D_\varepsilon$ is relatively ample.  Relative klt base point
freeness therefore makes a positive multiple of $L$ semiample over
$V$; the applicable statement is
\cite[Theorem~3.9.1]{BCHM}.  Its associated morphism contracts every
curve in every $f$-fibre.  A projective image of positive dimension
contains a curve, and a curve in such an image has a curve lift.
Consequently that morphism is constant on each connected $f$-fibre.
Normality and the usual factorization through a proper morphism with
connected fibres show that it factors through $V$.  Equivalently,
after a further multiple its semiample model is finite and
birational over $V$, hence is $V$.  Thus
\[
 L\sim_{\mathbb Q} f^*M
\]
for a $\mathbb Q$-Cartier divisor $M$ on $V$.

Since the flip is small on both sides, this identity of rational
line bundles transforms into
\[
 K_{U^+}+\Delta^+
   \sim_{\mathbb Q}tD_\varepsilon^++(f^+)^*M.
\]
The full adjoint is therefore ample over $V$.  The graph comparison
of Lemma~\ref{lem:graph} applies to these adjoints and their strict
boundaries and proves discrepancy improvement; in particular the
full positive pair is lc.  The same comparison gives the usual dlt
preservation.  More explicitly, every valuation of log discrepancy
zero on the positive side had discrepancy zero on the negative
side, and its centre was not contained in the bad locus.  Its generic
centre is consequently in the unchanged snc neighbourhood supplied
by the dlt property.  The characterization of dlt by such snc
neighbourhoods, or directly
\cite[Lemma~3.10.10(1)]{BCHM}, gives dlt on $U^+$.

For an elementary divisorial contraction the target is
$\mathbb Q$-factorial by the usual klt contraction theorem, so the
pushforward full adjoint is $\mathbb Q$-Cartier.  The same discrepancy
argument gives dlt preservation; see also
\cite[Remark~3.10.5 and Lemma~3.10.10(1)]{BCHM}.
\end{proof}

\begin{theorem}\label{thm:dlt-termination}
Every program of elementary negative birational steps starting from
a projective $\mathbb Q$-factorial dlt fourfold pair with rational boundary has finitely
many steps, independently of its successive ray choices.
\end{theorem}

\begin{proof}
Suppose such a program were infinite.  By
Lemma~\ref{lem:cycle-rank}, only finitely many of its steps lose a
prime divisor.  Discard that finite prefix, leaving a sequence of
elementary dlt flips.

We use special termination in its ordinary dlt form:
\cite[Theorem~4.2.1 and Remark~4.2.2]{FujinoSpecial}.  This theorem
assumes the log MMP for $\mathbb Q$-factorial dlt pairs in dimensions
at most one less than the dimension in question.  For rational
boundaries it needs only the rational-boundary forms of those
lower-dimensional results.  Here they are the classical
at-most-three-dimensional log MMP, including flip existence and
arbitrary termination, as in \cite{KollarEtAl1992}.  In particular
we do not invoke termination for four-dimensional lc pairs.  The
required fourfold special-termination conclusion is also stated
directly in \cite[Theorem~1.2]{FujinoFourfold}: eventually both the
flipping and the flipped loci are disjoint from the round-down of
the boundary.

After this further truncation, replace each boundary $\Delta_i$ by
$\Delta_i-\lfloor\Delta_i\rfloor$.  These pairs are klt.  The
removed divisors miss the surgery on both sides, so they have degree
zero on all curves of the two elementary contractions.  Thus the
same negative rays and the same two relative signs remain valid
for the decreased boundaries.  Relative ampleness is unchanged
by this relatively numerically trivial difference.  We have
obtained an infinite program from a projective klt fourfold pair with rational
boundary, contrary to Theorem~\ref{thm:klt-termination}.
\end{proof}

\subsection{Lifting every prescribed lc tail}

Lifting prescribed birational steps to finite dlt programs is standard;
see \cite[Lemma~3.1]{LazicMoragaTsakanikas} and
\cite[Lemma~2.6]{HanLiuZhuang2025}; compare the scaling construction in
\cite[Remark~2.9]{BirkarLC}.  We record the construction here, including
the nonemptiness of each bridge and its identification with the
prescribed positive model.

The following elementary convex observation is useful because a
crepant dlt lift has adjoint-zero curves over the model downstairs.

\begin{lemma}\label{lem:vertical-ray}
Let $(U,\Delta)$ be a projective lc pair, and let $h:U\to Z$ be a
projective morphism to a projective variety.  If $K_U+\Delta$ is not
nef over $Z$, there is a $(K_U+\Delta)$-negative extremal ray of the
absolute cone $\overline{\mathrm{NE}}(U)$ which is contracted by $h$.
Its contraction factors $h$.  If $h$ is birational, that contraction
is birational.
\end{lemma}

\begin{proof}
Choose ample divisors $A$ on $U$ and $H$ on $Z$.  The cone
\[
 F_h=\overline{\mathrm{NE}}(U)\cap(h^*H)^\perp
\]
is a closed face of the absolute cone.  A negative vertical curve
exists by the failure of relative nefness.  Hence the linear
functional $K_U+\Delta$ has negative minimum on the compact slice
$F_h\cap\{A\cdot\gamma=1\}$.  It has a negative extreme point
there, which generates an extremal ray of $F_h$ and therefore of
$\overline{\mathrm{NE}}(U)$.  The lc cone theorem represents this
negative ray by a curve.  Its zero intersection with $h^*H$ says
that it is vertical, since $H$ is ample.

The absolute contraction contracts only curves in this ray.
It follows that $h$ is constant on each of its connected fibres:
otherwise a positive-dimensional projective image would contain
a curve, which could be lifted to a curve in the fibre.  The
factorization through its normal target follows by the usual
proper connected-fibre factorization.  If $h$ is birational, a
fibre-type contraction would force the dimension of that target
to be smaller than $\dim Z=\dim U$, although it dominates $Z$.
This is impossible.
\end{proof}

\begin{proposition}\label{prop:lc-lift}
Every program in Definition~\ref{def:program} starting from a
projective lc fourfold pair over $\mathbb C$ with rational boundary has finitely many
birational steps.
\end{proposition}

\begin{proof}
Suppose the contrary.  Lemma~\ref{lem:cycle-rank} discards only
finitely many steps and leaves an infinite small tail
\[
 (X_i,B_i)\longrightarrow Z_i\longleftarrow(X_{i+1},B_{i+1}),
 \qquad i\ge i_0,
\]
in which both morphisms are small.  All choices in this tail are
kept fixed.

Take a projective crepant $\mathbb Q$-factorial dlt modification
\[
 p_{i_0}:(U_{i_0},\Delta_{i_0})\longrightarrow(X_{i_0},B_{i_0}).
\]
The dlt blowup theorem provides such a morphism, with
\[
 \Delta_{i_0}=(p_{i_0})_*^{-1}B_{i_0}
                   +\operatorname{Exc}(p_{i_0})_{\mathrm{red}},
 \qquad
 K_{U_{i_0}}+\Delta_{i_0}=p_{i_0}^*D_{i_0};
\]
see \cite[Theorem~10.4]{FujinoFoundations}.  In particular every
extracted prime has coefficient one and log discrepancy zero.

Inductively suppose that $p_i:(U_i,\Delta_i)\to(X_i,B_i)$ is such
a crepant dlt lift.  Over $Z_i$ run elementary dlt MMP steps for
$K_{U_i}+\Delta_i$ until it becomes relatively nef.  At any non-nef
intermediate stage,
Lemma~\ref{lem:vertical-ray} gives a negative \emph{absolute}
extremal ray inside the vertical face.  The face need not be
entirely negative.  The elementary step exists by
Lemma~\ref{lem:dlt-elementary}, its contraction factors the map
to $Z_i$, and it is birational because that map is birational.
Every total space is projective and $\mathbb Q$-factorial dlt.
An infinite choice of such relative steps would consequently be
an infinite absolute elementary dlt program, excluded by
Theorem~\ref{thm:dlt-termination}.  The chosen relative program
therefore reaches a nef end, say $(U_i',\Delta_i')$.

This finite bridge is nonempty.  Choose a curve $C$ contracted by
$X_i\to Z_i$.  It has $D_i\cdot C<0$.  Projectivity of $p_i$
provides a curve $\widetilde C$ dominating $C$: take a component
of the inverse image dominating $C$ and cut by sufficiently many
general ample divisors.  Crepancy gives
\[
 (K_{U_i}+\Delta_i)\cdot\widetilde C
     =\deg(\widetilde C/C)D_i\cdot C<0.
\]
Thus the starting lift was not nef over $Z_i$.

We identify the end with a crepant lift of the \emph{prescribed}
positive model $X_{i+1}$.  On a common smooth model of $U_i$,
$U_i'$, and $X_{i+1}$, the pullback of the starting adjoint can be
written in two ways.  The relative MMP writes it as the pullback
of $K_{U_i'}+\Delta_i'$ plus an effective divisor exceptional over
$U_i'$.  Crepancy of $p_i$ and Lemma~\ref{lem:graph} write it as
the pullback of $D_{i+1}$ plus an effective divisor exceptional over
$X_{i+1}$.  One end is nef over $Z_i$ and the other is ample there.
Lemma~\ref{lem:two-ends} gives a projective birational morphism
\[
 p_{i+1}:U_i'\longrightarrow X_{i+1},\qquad
 K_{U_i'}+\Delta_i'=p_{i+1}^*D_{i+1}.
\]
Every prime exceptional for $p_{i+1}$ is the transform of a prime
already exceptional for $p_i$.  Indeed the downstairs map is
small, every downstairs prime persists, and the elementary
program upstairs extracts no primes.  Each remaining extracted
prime therefore still has coefficient one.  Setting
$(U_{i+1},\Delta_{i+1})=(U_i',\Delta_i')$ supplies the next
crepant dlt lift without any additional modification.

Concatenating these finite nonempty bridges produces an infinite
absolute elementary dlt program from $(U_{i_0},\Delta_{i_0})$.
This again contradicts Theorem~\ref{thm:dlt-termination}.  The
argument has followed every prescribed downstairs step; it has
not replaced their ray choices by a program with scaling.
\end{proof}

\subsection{Continuation through mixed steps}

It remains to construct an output whenever the chosen negative
contraction is birational.  The relative lc existence theorem supplies
a good model over its base; the graph comparison will show that its
ample model gives a permitted step.

\begin{proposition}\label{prop:continuation}
Let $(X,B)$ be a projective lc fourfold pair over $\mathbb C$ with rational boundary,
with $D=K_X+B$ $\mathbb Q$-Cartier.  Every $D$-negative extremal
ray has a projective contraction with normal projective target,
connected fibres, relative Picard number one, and relatively
ample $-D$.  If this contraction is birational, it has an output
in the precise convention of Definition~\ref{def:program}:
the positive map is small, possibly an isomorphism, and its
actual adjoint with the strict boundary is relatively ample
and lc.  If it is of fibre type, it is a Mori fibre-space
endpoint in that convention.
\end{proposition}

\begin{proof}
The lc cone and contraction theorem applies without
$\mathbb Q$-factoriality; see
\cite[Theorem~1.1]{FujinoFoundations}.  It gives, for the chosen
negative extremal ray, the stated contraction $f:X\to Z$.
Its target is projective.  Relative Picard number one and the
negative degree of $D$ give relative ampleness of $-D$.
The fibre-type case already has precisely the required endpoint
properties, so assume $f$ birational.

Choose a sufficiently divisible integer $r>1$ and an ample
Cartier divisor $H$ on $Z$ such that
\[
 \mathcal O_X(-rD+f^*H)
\]
is generated by global sections.  This is possible by relative
ampleness and an ample base twist.  A general member $G$ gives
an effective rational $\mathbb Q$-Cartier divisor $A=G/r$ with
\[
 D+A\sim_{\mathbb Q}0/Z,
 \qquad (X,B+A)\ \text{lc}.
\]
For the latter assertion, take a log resolution of $(X,B)$ and
choose $G$ generally in the free system.  Its pullback has no
exceptional component and is transverse to the finitely many
SNC strata on the resolution.  Adding its coefficient $1/r$
keeps every coefficient at most one in the crepant log-smooth
calculation.  This proves lc, including when the original pair
has coefficient-one components.

We apply \cite[Theorem~1.1(1)--(2)]{BirkarLC} in its published
rational complemented-pair form.  Its hypotheses are that
$(X/Z,B+A)$ is lc, $B,A$ are effective rational divisors, $A$ is
$\mathbb Q$-Cartier, the structural morphism is projective and
surjective, and $K_X+B+A\sim_{\mathbb Q}0/Z$.  It does not assume
$\mathbb Q$-factoriality or ampleness of $A$.  It gives a log
minimal model or a Mori fibre space for $(X/Z,B)$, and the
adjoint on a nef log minimal model is semiample over $Z$.
A Mori fibre space over $Z$ is impossible here: its intermediate
base would have dimension less than four and would dominate the
four-dimensional birational base $Z$.  We therefore obtain a
good log minimal model $(M,B_M)$ over $Z$.

Let $q:M\to V$ be its relative semiample contraction.  Since
$M\to Z$ is birational, this is a birational contraction, and
$V\to Z$ is projective and birational.  There is an ample-over-$Z$
$\mathbb Q$-Cartier divisor $L$ on $V$ with
$K_M+B_M\sim_{\mathbb Q}q^*L$.  Put $B_V=q_*B_M$.  Pushing the
linear equivalence forward shows that the actual divisor
\[
 D_V=K_V+B_V
\]
is $\mathbb Q$-Cartier and $\mathbb Q$-linearly equivalent to $L$.
Moreover
\[
 K_M+B_M=q^*D_V.
\]
Indeed the difference is $q$-exceptional and relatively
numerically trivial, so negativity in both signs makes it zero.
Thus the ample divisor on $V$ is its actual adjoint, not just a
numerical substitute.

The log-minimal-model discrepancy comparison gives, on a smooth
common resolution $a:W\to X$, $b:W\to M$,
\[
 a^*D=b^*(K_M+B_M)+E,
 \qquad E\ge0,
 \qquad E\ \text{exceptional over }M;
\]
see \cite[Remark~2.6]{BirkarLC}.  This also explains possible
extractions in the definition of a log minimal model.  A prime
extracted on $M$ is given coefficient one; its coefficient in
$E$ is then $-a(v;X,B)\le0$.  Effectivity forces this coefficient
to vanish, so that extraction is crepant.  After composition
with $q$, the same $E$ is effective and exceptional over $V$.

We prove that $V\to Z$ is small, although $X\to Z$ need not be.
Let $G_0$ be the normalized graph of $X\dashrightarrow V$, with
projections $u:G_0\to X$ and $v:G_0\to V$, and set
\[
 H_0=u^*D-v^*D_V.
\]
The pullback of $H_0$ to a smooth model dominating $W$ and $G_0$
is $E$ pulled back.  Pushing down gives
\[
 H_0\ge0,\qquad v_*H_0=0.
\]
Its negative is ample over $Z$: it is the sum of the pullbacks
of $-D$ and $D_V$ from the two factors, and the graph maps
finitely into their fibre product.  As in the support argument
of Lemma~\ref{lem:graph}, every graph fibre over a point where
one of the two maps to $Z$ is nonisomorphic is contained in
$\operatorname{Supp}H_0$.  Indeed that fibre is connected and
positive-dimensional, every closed point lies on a contracted
curve, and anti-ampleness contradicts effectivity on a curve
not contained in the support.

If a prime divisor on $V$ were exceptional over $Z$, the preceding
support property would force a component of $H_0$ to dominate
it.  This contradicts $v_*H_0=0$.  Hence $V\to Z$ is small.
Every prime on $V$ consequently corresponds to a prime on $X$
through their common codimension-one open over $Z$.  Thus
$B_V$ is exactly the strict transform of $f_*B$; there are no
additional boundary primes on $V$.  The effective discrepancy
difference gives
\[
 a(\xi;X,B)\le a(\xi;V,B_V)
\]
for every divisorial valuation, so $(V,B_V)$ is lc.  The diagram
$X\to Z\leftarrow V$ is the required step, including a mixed
step when $f$ loses a divisor and $V\to Z$ is a nontrivial small
modification.  No positive-side relative Picard-number assertion
has been used.
\end{proof}

\begin{remark}
The existence input in Proposition~\ref{prop:continuation} is
Theorem~1.1 of \cite{BirkarLC}, not merely its Corollary~1.2,
whose literal formulation begins with a small contraction.
The effective supplement and the graph argument supply the
additional mixed-step scope needed here.  The ACC-dependent
generalization in that paper is unnecessary for this rational
complemented application.
\end{remark}
